\chapter[Sermon 14. That the Doctrine of Piety Is So Fully Set Forth to the Church, That There Needs No New Revelations. And of the Most Large Promises of Christ Made unto the Church.]{That the Doctrine of Piety Is So Fully Set Forth to the Church, That There Needs No New Revelations. And of the Most Large Promises of Christ Made unto the Church.}
\label{sermon:014}
\markright{Sermon 14. That the Doctrine of Piety Is So Fully Set Forth to ...}

And to you I say, and to others who are of Thyatira: Whoever does not have this doctrine, and has not known the devices of Satan, as they say, I will place no other burden upon you than what they already have. Hold fast until I come. And whoever overcomes and keeps my works until the end, to him I will give power over nations, and he shall rule them with a rod of iron. And as a potter breaks the vessels, so will he break them. Even as I received from my father, so I will give it to him—the morning star. Let him who has ears hear what the Spirit says to the congregations.

\index{Papacy}He now speaks to the Cataphrygians and also to the faithful of the church of Thyatira who rightly believed in Christ, and heals their diseases. This demonstrates the unspeakable mercy of God, which does not cease to speak to those still entangled with heresy and to heal their pestilential diseases. He admonishes all men that they should not look for new revelations, but rather know that God through Christ and his Apostles has set forth a most perfect doctrine, to which he will add nothing. Therefore, they should keep fast in memory the things they had already learned and in which they were now being exercised. For the Cataphrygians, also called Montanists, boasted of a new comforter and a new revelation, as though all things had not been fully set forth by the Apostles, but that many things were still to be revealed. As also the maintainers of the Popish church stubbornly affirm today. And just as the Cataphrygians covered their trivialities under the pretense of the Holy Spirit, so too do the Papists cloak the vain constitutions of men and set them forth under a false color of the Holy Spirit, as though the Lord spoke of their decrees when he said, “I have yet many things to say unto you, which now you cannot bear.” Nevertheless, the faithful people of Thyatira, who did not have the doctrine of Jezebel but rather detested it as doubtful, said that the Devil was a certain dependence [?] and had a thousand crafts, which could also transform him into an angel of light. And that they were but simple men, who, being ignorant of his wonderful crafts and subtleties, did not know what they should chiefly follow, while false prophets also make their boast of the Holy Spirit, shine in miracles, and avouch their doctrine to be true. You will find this even today, when people will say, “I am a plain, simple man, and do not know to which side I should cleave, since the doctors of both sides affirm with great cost that they have the truth on their side, and therefore some will say, they shall agree better before I will believe any of them.” \&c.

The Lord, therefore, answering both parties, shows what they should do. To you, he says, who follow the doctrine of Jezebel, I say also to the rest of the Thyatirans, who do not follow Jezebel’s doctrine, yet nevertheless complain in such disagreements and wondrous schemes of the devil, that they do not see what is best: To you all I say, if you are simple in deed, as you pretend, if you will with all your heart embrace the truth, give yourselves to the simplicity of the apostolic faith, cleaving fast to such things as you have once learned from the Apostles, neither seeking nor receiving any new religions, or additions, constitutions, or anything else moreover, than that which you have learned from the Apostles. For these things which you have received are sufficient to obtain salvation.

And these words of the Lord must be weighed more diligently, so that we may perceive the great benefit in them: that is, I will lay upon you no other weight or burden besides what you already have. The Lord affirms that he will add nothing more to the evangelical doctrine set forth by the Apostles, as to that which is most perfect. Certainly, if the doctrine of Moses were so perfect that the Lord himself prohibited any man from adding or taking away anything from it, but only doing what was commanded, as we read in the fourth and twelfth chapters of Deuteronomy. Who would doubt that the doctrine of Christ, the Son of God, would be lacking in anything? He therefore now affirms that he will lay nothing upon them more than he had already laid, and which they were bearing at that time.

A burden in the sermons of the Prophets is taken as doctrine of grave and weighty matters. The Apostles also call the law a yoke and burden. Wherefore, when the Lord says that he will not lay any other burden upon the church, he means that he will not reveal any other doctrine, nor further charge them with other rites or ceremonies than those he had already ordained and imposed. These words of Christ accord very well with what is read in the apostles’ epistle, the Synodical Act of xv. For by the common consent of the congregation, and after the mind of the Holy Spirit, they say they will impose nothing moreover upon the church than such things as they had already received of Paul, and a few things that they added for a declaration of the same. Whereupon Paul said to the Galatians: If an angel from heaven preach unto you another gospel besides that which is preached, let him be accursed.

\index{Arethas of Caesarea}\index{Judgment, Last}What then? Hold fast, namely that which you have received, suffering it not to be plucked out of your hands: Hold fast I say, with tooth and nail, until I come; that is to say, unto the last judgment. Therefore he testifies expressly that this doctrine shall be perpetual and unchangeable, and therefore to be kept most steadfastly by all of us, and not to be shaken from, though all the world cry out and persuade the contrary. Arethas, Bishop of Caesarea, required of them nothing else, he says, but that they would keep safely the godly pledge of faith until his coming. This if we shall do, we may easily avoid the crafts of the devil and deceivable clouds. For whatsoever they shall bring forth, whatsoever they shall forge and refine, or disguise with the counterfeited color of the Holy Spirit, we shall always have recourse to the simple doctrine of Christ set forth by the apostles, wherein alone we shall rest, rejecting all things that shall not accord with the same.

\index{Rome}And this wholesome doctrine of Christ confounds all traditions, and subverts all constitutions made since the time of the Apostles. The godly may always object this saying of Christ to the traditioners. I will lay no other burden upon you, besides that you have: That same hold fast until the last judgment. They shall allege that same also, that the Apostles deny that they will add anything more. Acts 15. Christ spake this in the time of John, in the year of our Lord 87. Therefore whatever laws, traditions, decrees have been made since that time, we know they were not imposed by Christ, which saith so expressly that he will lay no other burden on the faithful. Where then become the decrees and constitutions of worshipping images in the church, for the consecration and celebrating of masses? What shall we say to the decretals of the Bishop of Rome? They are all overthrown and stricken down as it were with a thunderbolt, by this only sentence of Christ. I will impose no other burden than that you have, keep that until the judgment. Behold, he saith, unto the judgment, lest any should imagine in the meantime that another thing had pleased the Holy Spirit. Let us therefore persevere in the same.

Hereunto he adds, as was his custom, most ample promises, that through hope of so great rewards he might draw them from errors and join them to the true religion. And just as in the former epistles he has said, “He who overcomes,” so here he repeats the same, admonishing us not to sleep but to watch and fight valiantly. He overcomes who keeps the works of Christ until the end. The works of Christ, by a subtle opposition, are set against the inventions and works of men.

The works of Christ signify both doctrine and faith, and whatever good works flow from them, the service of worshiping God, and the observance of God’s word. For in the twenty-eighth chapter of Saint Matthew, the Lord says to his disciples, “Teach them to keep those things which I have commanded you.” He speaks with emphasis, “those things which I have commanded you,” not such as you shall have invented of your own brain. For the Lord quotes from the prophet in the fifteenth chapter of the same Saint Matthew, saying, “In vain do they worship me, teaching the doctrines of men.” Therefore these works have no promise; but the works of Christ, which he himself has ordained, and which are done by his Spirit and by true faith, while we forsake our errors and cleave to the truth, they have a most ample promise.

And promises two notable things. The first: just as my father has promised me victory, and fulfilled it, that I overcome all my enemies, and triumphed over them, shattering them like clay vessels without any difficulty, so will I give unto you also power and victory against all the ungodly. And that same promise shall at the last be fully accomplished in the final judgment, in which all the enemies of godliness shall be cast under the feet of Christ, as it is declared in the Psalms, especially in the second and one hundredth Psalm. And in this world also Christ affirms that his servants shall spiritually rule over his enemies, just as Christ, although he was tormented and died, nevertheless overcame his enemies. The holy and ecclesiastical histories bear witness of these things sufficiently.

\index{Resurrection}\index{Daniel (Book of)}The latter: I will give him the morning star. And he understood the knowledge of Christ increasing daily more and more, and so even Christ himself, in like case as the day in the rising of the morning star waxes brighter and brighter. In this sense the Apostle Saint Peter is said to have used this allegory in the first chapter of his Second Epistle, or at the least he promised a clarity most bright. For Daniel says how the faithful in the resurrection shall shine like the firmament. The which thing also the Lord Christ mentions in Matthew 13. And the Apostle alluding hereunto said that one star was brighter than another: so likewise in the resurrection one shall be made brighter than another. These promises are most great, neither can I think that any greater can be given us. God grant us grace that we may be made partakers of so great things.

Finally, he applies this epistle to all the churches and all ages of the world. Concerning which, since we have spoken of it more than once, there is no need for me to be tedious to anyone by repeating it again. To the Lord our God be praise and glory.
